% ============================================================
%  ch6_experimental_evaluation.tex
%  Chapter 6: Experimental Evaluation
%  Master's Thesis -- SQL MATCH_RECOGNIZE on Pandas DataFrames
% ============================================================
\chapter{Experimental Evaluation}
\label{chap:evaluation}

This chapter evaluates the proposed \texttt{MATCH\_RECOGNIZE}
implementation across five dimensions: output correctness on selected
Trino-aligned reference queries, coverage of the supported SQL:2016
row-pattern subset, row-pattern type coverage, performance and
scalability, and developer usability.
Section~\ref{sec:eval-methodology}
describes the evaluation methodology, including the hardware
environment and the metrics used.
Section~\ref{sec:eval-dataset} presents the benchmark dataset and
test patterns.
Section~\ref{sec:eval-correctness} reports correctness results.
Section~\ref{sec:eval-performance} presents performance results.
Section~\ref{sec:eval-comparison} compares the system with existing
tools.

% ----------------------------------------------------------------
\section{Evaluation Methodology}
\label{sec:eval-methodology}
% ----------------------------------------------------------------

This chapter evaluates the proposed \texttt{MATCH\_RECOGNIZE} implementation
through five complementary criteria.

\begin{enumerate}
  \item \textbf{Correctness.}
    Results are compared row-by-row with Trino across selected
    pattern types to check matched rows, pattern alignment, and
    measure calculations for the supported feature subset.

  \item \textbf{SQL:2016 feature coverage.}
    Support for key \texttt{MATCH\_RECOGNIZE} clauses is verified
    against the implemented subset, including \texttt{PARTITION~BY},
    \texttt{ORDER~BY},
    \texttt{PATTERN}, \texttt{DEFINE}, \texttt{MEASURES},
    \texttt{SUBSET}, output modes (\texttt{ONE~ROW~PER~MATCH} /
    \texttt{ALL~ROWS~PER~MATCH}), after-match skip policies, and
    navigation
    functions~\cite{iso2016,ISO2021,monierashraf2025rowmatch}.

  \item \textbf{Pattern-type coverage.}
    A variety of row-pattern types are exercised, including simple
    sequences, quantified repetitions
    (\texttt{*}, \texttt{+}, \texttt{?}, \texttt{\{n,m\}}),
    ascending/descending trends, recovery patterns, alternation,
    exclusion blocks (\texttt{\{-\ldots-\}}), anchored patterns,
    and \texttt{PERMUTE} constructs.

  \item \textbf{Performance and scalability.}
    Execution time is measured across increasing DataFrame sizes
    and pattern complexities to characterize how the engine scales
    with data volume and pattern structure.

  \item \textbf{Usability (developer effort).}
    The lines of code required to express equivalent logic as
    idiomatic Pandas serve as a baseline.  Hand-written Pandas
    solutions typically require substantially more code
    (e.g.\ chained \texttt{groupby}/\texttt{shift}/joins) than a
    single \texttt{MATCH\_RECOGNIZE} statement, so reducing LOC
    is a core goal of the RPR
    approach~\cite{iso2016,ISO2021,melton2017sqlrpr,Python_UDF}.
\end{enumerate}

% ................................
\subsection{Environment and Configuration}
\label{subsec:eval-environment}
% ................................

The experiments were performed on a workstation running Ubuntu
24.04.2 LTS with an Intel Core i7-11800H (16 cores at 2.30\,GHz)
and 16\,GB of DDR4 RAM at 3200\,MHz.

The software environment comprised:
\begin{itemize}
  \item \textbf{Baseline SQL engine:} Trino~473.
  \item \textbf{Runtimes:} Java~11.0.25 LTS and Python~3.12.7.
  \item \textbf{Core Python libraries:} \texttt{pandas}~2.2.3,
    \texttt{numpy}~1.26.4,
    \texttt{antlr4-python3-runtime}~4.13.2,
    \texttt{antlr4-tools}~0.2.1,
    \texttt{pytest}~8.3.4,
    Jupyter,
    \texttt{unittest},
    \texttt{matplotlib}~3.10.0,
    \texttt{plotly}~5.24.1, and
    \texttt{seaborn}~0.13.2.
  \item \textbf{Standard library modules used:}
    \texttt{typing}, \texttt{dataclasses}, \texttt{re},
    \texttt{operator}, \texttt{math}, \texttt{itertools},
    \texttt{collections}, \texttt{enum}, \texttt{time},
    \texttt{sys}, \texttt{os}, \texttt{logging}, and
    \texttt{traceback}.
\end{itemize}

The \texttt{MATCH\_RECOGNIZE} implementation uses ANTLR4 for
parsing Trino SQL, \texttt{pandas} for DataFrame management, and
custom AST/NFA/DFA modules for pattern detection (see
Chapter~\ref{chap:methodology}).

% ................................
\subsection{Evaluation Metrics}
\label{subsec:eval-metrics}
% ................................

Each test case is assessed on the five dimensions defined above.

\paragraph{Correctness.}
The result DataFrame produced by the implementation is compared
with the result set returned by Trino using an exact row-level
equality check---same rows, same column values, and same order
where order is deterministic.

\paragraph{SQL:2016 feature coverage.}
Each clause and sub-feature of \texttt{MATCH\_RECOGNIZE} is marked
as \emph{supported}, \emph{partially supported}, or
\emph{unsupported} based on whether the implementation supports the
feature and whether the exercised behavior matches the Trino-aligned
reference tests.

\paragraph{Pattern-type coverage.}
Test cases are grouped by row-pattern type
(sequence, repetition, trend, recovery, alternation, exclusion,
anchor, \texttt{PERMUTE}).  A pattern type is considered covered
if at least one test case exercises it and produces correct output.

\paragraph{Performance and scalability.}
Execution time (wall-clock seconds, averaged over three runs) is
recorded for increasing DataFrame sizes and for patterns of
varying complexity.  Measurements report mean execution time and
are used to characterize how the engine scales with data volume
and pattern structure.

\paragraph{Usability.}
The number of non-blank, non-comment lines in the equivalent
idiomatic Pandas solution is recorded and compared with the
single-statement \texttt{MATCH\_RECOGNIZE} query, providing a
concrete measure of developer effort reduction.

% ----------------------------------------------------------------
\section{Dataset and Test Patterns}
\label{sec:eval-dataset}
% ----------------------------------------------------------------

% ................................
\subsection{Dataset}
\label{subsec:eval-dataset-orders}
% ................................

An \texttt{orders} table is used to demonstrate and validate
pattern-matching queries.  The dataset comprises three
attributes~\cite{trino2021row}: \texttt{customer\_id}
(\texttt{VARCHAR}), which identifies the customer (two distinct
values: \texttt{cust\_1} and \texttt{cust\_2});
\texttt{order\_date} (\texttt{DATE}), the date the order was placed
(ranging from 2020-05-11 to 2020-05-18); and \texttt{price}
(\texttt{INTEGER}), the purchase amount on that date.
Table~\ref{tab:orders} shows the eight records used throughout the
experiments.

\begin{table}[H]
  \centering
  \caption{Sample \texttt{orders} dataset
    (\texttt{customer\_id}, \texttt{order\_date}, \texttt{price})}
  \label{tab:orders}
  \begin{tabularx}{\textwidth}{|>{\centering\arraybackslash}X|>{\centering\arraybackslash}X|>{\centering\arraybackslash}X|}
    \hline
    \textbf{customer\_id} & \textbf{order\_date} & \textbf{price} \\
    \hline
    cust\_1 & 2020-05-11 & 100 \\
    cust\_1 & 2020-05-12 & 200 \\
    cust\_2 & 2020-05-13 &   8 \\
    cust\_1 & 2020-05-14 & 100 \\
    cust\_2 & 2020-05-15 &   4 \\
    cust\_1 & 2020-05-16 &  50 \\
    cust\_1 & 2020-05-17 & 100 \\
    cust\_2 & 2020-05-18 &   6 \\
    \hline
  \end{tabularx}
\end{table}

Each customer appears multiple times in chronological order:
\texttt{cust\_1} spans five dates and \texttt{cust\_2} spans three.
This structure enables focused validation of partitioning, ordering,
and pattern detection (e.g.\ sequences of rising or falling prices per
customer).  In the experiments, data are loaded via a
\texttt{VALUES} clause in Trino and checked against the Pandas
implementation.

\begin{lstlisting}[basicstyle=\ttfamily\small,
  backgroundcolor=\color{gray!10},
  caption={Creating the in-memory table in Trino}]
CREATE TABLE memory.default.orders (
    customer_id VARCHAR,
    order_date  DATE,
    price       INTEGER
);
\end{lstlisting}

\begin{lstlisting}[basicstyle=\ttfamily\small,
  backgroundcolor=\color{gray!10},
  caption={Inserting data into the table}]
INSERT INTO memory.default.orders VALUES
    ('cust_1', DATE '2020-05-11', 100),
    ('cust_1', DATE '2020-05-12', 200),
    ('cust_2', DATE '2020-05-13',   8),
    ('cust_1', DATE '2020-05-14', 100),
    ('cust_2', DATE '2020-05-15',   4),
    ('cust_1', DATE '2020-05-16',  50),
    ('cust_1', DATE '2020-05-17', 100),
    ('cust_2', DATE '2020-05-18',   6);
\end{lstlisting}

% ----------------------------------------------------------------
\section{Correctness Results}
\label{sec:eval-correctness}
% ----------------------------------------------------------------

This section demonstrates output correctness by comparing results from
three independent approaches---the proposed engine, Trino~473, and a
hand-written Pandas baseline---on the same SQL \texttt{MATCH\_RECOGNIZE}
query and dataset.

% ................................
\subsection{Experimental Setup}
\label{subsec:correctness-setup}
% ................................

The correctness experiments use the same hardware and software
configuration described in Section~\ref{subsec:eval-environment}.
For each reference query, the input data are loaded into both Trino
and a Pandas DataFrame with equivalent schemas.  The result sets are
then compared by row order and measure values after normalizing
case-only column-name differences between SQL and Pandas outputs.

% ................................
\subsection{Use Case: Stock Price Analysis}
\label{subsec:eval-stock-price}
% ................................

The \texttt{stock\_price} table is used to demonstrate and validate
pattern-matching queries with the \texttt{MATCH\_RECOGNIZE} clause.
The table includes three columns: \texttt{company}, \texttt{price\_date},
and \texttt{price}, representing a history of daily stock prices for
market analysis.  Here, \texttt{company} specifies the company name,
\texttt{price\_date} is the trading date, and \texttt{price} records
the stock's value on that day~\cite{petkovic2022rpr}.

This example detects bullish momentum by identifying a three-row
window in which the second and third rows increase relative to their
immediate predecessors.  The query partitions records by company
(\texttt{PARTITION BY company}), orders them chronologically
(\texttt{ORDER BY price\_date}), and applies \texttt{PATTERN (A B C)}
with \texttt{DEFINE} predicates on \texttt{B} and \texttt{C}.  The
variable \texttt{A} has no explicit \texttt{DEFINE} predicate and
therefore uses the standard implicit always-true behavior.  Key
attributes---starting price, ending price, and match ID---are produced
in the \texttt{MEASURES} clause.  A structured \texttt{pandas}
DataFrame serves as input to \texttt{MATCH\_RECOGNIZE}, enabling
automated detection of upward trends across firms for the supported
row-pattern subset.
Additional use cases are available in the project
repository~\cite{monierashraf2025rowmatch}.

\begin{table}[H]
  \centering
  \caption[Stock price table]{The tabular form of the \texttt{stock\_price} table}
  \label{tab:stock_price}
  \begin{tabular}{@{}lll@{}}
    \toprule
    \texttt{Company} & \texttt{Price\_date} & \texttt{Price} \\
    \midrule
    A & 2020-10-01 & 50 \\
    B & 2020-10-01 & 89 \\
    A & 2020-10-02 & 36 \\
    B & 2020-10-02 & 24 \\
    A & 2020-10-03 & 39 \\
    B & 2020-10-03 & 37 \\
    A & 2020-10-04 & 42 \\
    B & 2020-10-04 & 63 \\
    A & 2020-10-05 & 30 \\
    B & 2020-10-05 & 65 \\
    A & 2020-10-06 & 47 \\
    B & 2020-10-06 & 56 \\
    A & 2020-10-07 & 71 \\
    B & 2020-10-07 & 50 \\
    A & 2020-10-08 & 80 \\
    B & 2020-10-08 & 54 \\
    A & 2020-10-09 & 75 \\
    B & 2020-10-09 & 30 \\
    A & 2020-10-10 & 63 \\
    B & 2020-10-10 & 32 \\
    \bottomrule
  \end{tabular}
\end{table}

\noindent\textbf{Load the dataset as a Pandas DataFrame.}

\begin{lstlisting}[language=Python,
  basicstyle=\ttfamily\small,
  backgroundcolor=\color{gray!10},
  caption={[Python baseline dataset]Python baseline dataset.},
  label={lst:dataset}]
import pandas as pd
data = {
    "Company": [
        "A","B","A","B","A","B","A","B","A","B",
        "A","B","A","B","A","B","A","B","A","B"
    ],
    "Price_date": [
        "2020-10-01","2020-10-01","2020-10-02","2020-10-02","2020-10-03",
        "2020-10-03","2020-10-04","2020-10-04","2020-10-05","2020-10-05",
        "2020-10-06","2020-10-06","2020-10-07","2020-10-07","2020-10-08",
        "2020-10-08","2020-10-09","2020-10-09","2020-10-10","2020-10-10"
    ],
    "Price": [
        50,89,36,24,39,37,42,63,30,65,
        47,56,71,50,80,54,75,30,63,32
    ]
}
df = pd.DataFrame(data)
df["Price_date"] = pd.to_datetime(df["Price_date"])
\end{lstlisting}

\noindent\textbf{Using the proposed \texttt{MATCH\_RECOGNIZE}.}

\begin{lstlisting}[language=Python,
  basicstyle=\ttfamily\small,
  backgroundcolor=\color{gray!10},
  caption={[Using the proposed MATCH\_RECOGNIZE]Using the proposed \texttt{MATCH\_RECOGNIZE}.}]
from src.executor.match_recognize import match_recognize

query = """
SELECT *
FROM stock_price
MATCH_RECOGNIZE (
  PARTITION BY Company
  ORDER BY Price_date
  MEASURES
    FIRST(A.Price) AS start_price,
    LAST(A.Price)  AS end_price,
    MATCH_NUMBER() AS match_id
  PATTERN (A B C)
  DEFINE
    B AS B.Price > PREV(B.Price),
    C AS C.Price > PREV(C.Price)
); """
result = match_recognize(query, df)
print(result)
\end{lstlisting}

\noindent\textbf{Results (proposed engine):}
\begin{lstlisting}[basicstyle=\ttfamily\small, backgroundcolor=\color{gray!10}]
  Company  start_price  end_price  match_id
0       A           36         36         1
1       A           30         30         2
2       B           24         24         1
\end{lstlisting}

\noindent\textbf{Using Trino and Oracle as Baselines.}

In the second step, the same pattern-matching query was applied using the
same data in the Trino environment to validate the correctness of the proposed
approach.  The stock price history was created as an in-memory table
using \texttt{CREATE TABLE} and populated through \texttt{INSERT INTO},
ensuring that the schema and data (company, date, and price) remained
consistent with the Pandas DataFrame used earlier.  The
\texttt{MATCH\_RECOGNIZE} query was then executed in Trino to detect
the same three-row increasing-window pattern per company.  In
addition, the same query was validated in Oracle Database via
\texttt{MATCH\_RECOGNIZE}~\cite{petkovic2022rpr,monierashraf2025oracle},
using an equivalent \texttt{stock\_price} table with the same schema and
rows; the Oracle results matched the expected behavior, strengthening
cross-system validation under standard \texttt{MATCH\_RECOGNIZE}
semantics~\cite{monierashraf2025oracle}.

\begin{lstlisting}[basicstyle=\ttfamily\small,
  backgroundcolor=\color{gray!10},
  caption={[Create table in Trino]Creating the in-memory table in Trino.}]
CREATE TABLE memory.default.stock_price (
    Company    VARCHAR,
    Price_date DATE,
    Price      INTEGER
);
\end{lstlisting}

\begin{lstlisting}[basicstyle=\ttfamily\small,
  backgroundcolor=\color{gray!10},
  caption={[Insert rows in Trino]Inserting data into the table.}]
INSERT INTO memory.default.stock_price (Company, Price_date, Price) VALUES
('A', DATE '2020-10-01', 50),
('B', DATE '2020-10-01', 89),
('A', DATE '2020-10-02', 36),
('B', DATE '2020-10-02', 24),
...
('A', DATE '2020-10-10', 63),
('B', DATE '2020-10-10', 32);
\end{lstlisting}

\begin{lstlisting}[basicstyle=\ttfamily\small,
  backgroundcolor=\color{gray!10},
  caption={[Execute MATCH\_RECOGNIZE in Trino]Executing the same \texttt{MATCH\_RECOGNIZE} query in Trino.}]
SELECT *
FROM memory.default.stock_price
MATCH_RECOGNIZE (
  PARTITION BY company
  ORDER BY price_date
  MEASURES
    FIRST(A.price) AS start_price,
    LAST(A.price)  AS end_price,
    MATCH_NUMBER() AS match_id
  PATTERN (A B C)
  DEFINE
    B AS B.price > PREV(B.price),
    C AS C.price > PREV(C.price)
);
\end{lstlisting}

\noindent\textbf{Results (Trino):}
\begin{lstlisting}[basicstyle=\ttfamily\small, backgroundcolor=\color{gray!10}]
 company | start_price | end_price | match_id
---------+-------------+-----------+----------
 A       |          36 |        36 |        1
 A       |          30 |        30 |        2
 B       |          24 |        24 |        1
\end{lstlisting}

\noindent\textbf{Python DataFrame Baseline.}
As a Python baseline, the same pattern logic was implemented directly in
\texttt{pandas}, without using \texttt{MATCH\_RECOGNIZE}.  The input
DataFrame (Listing~\ref{lst:dataset}) is converted to \texttt{datetime}
and sorted by company and trading date.  For each company, the
baseline computes a Boolean signal
\texttt{inc\,:=\,diff(Price)\,>\,0} and marks candidate
starts where the next two differences are positive, corresponding to
the SQL pattern rows $A,B,C$ with $B>A$ and $C>B$.  Candidates are
processed sequentially with a
non-overlapping policy, advancing
\texttt{next\_start\_pos\,=\,i\,+\,3} to emulate
\texttt{AFTER MATCH SKIP PAST LAST ROW}.  For each accepted start, the
implementation outputs \texttt{company}, \texttt{match\_id}, and
measures \texttt{start\_price} and \texttt{end\_price} from the start
row~(A), returning the resulting matches as a
\texttt{DataFrame}~\cite{Python_UDF}.

\begin{lstlisting}[language=Python,
  basicstyle=\ttfamily\small,
  backgroundcolor=\color{gray!10},
  caption={[Equivalent logic in pandas]Equivalent pattern matching in Python/pandas.}]
import pandas as pd

def detect_two_consecutive_increases(
    df: pd.DataFrame,
    company_col: str = "Company",
    date_col:    str = "Price_date",
    price_col:   str = "Price",
) -> pd.DataFrame:
    gdf = df.copy()
    gdf[date_col] = pd.to_datetime(gdf[date_col])
    gdf = gdf.sort_values([company_col, date_col]).reset_index(drop=True)
    out_rows = []
    for company, g in gdf.groupby(company_col, sort=False):
        g = g.reset_index(drop=True)
        inc = g[price_col].diff() > 0
        start_mask = (inc.shift(-1).fillna(False)
                      & inc.shift(-2).fillna(False))
        candidates = [int(i) for i, ok in enumerate(start_mask)
                      if ok and i + 2 < len(g)]
        match_id, next_start_pos = 0, 0
        for i in candidates:
            if i < next_start_pos:
                continue
            match_id += 1
            out_rows.append({
                "company":     company,
                "start_price": g.loc[i, price_col],
                "end_price":   g.loc[i, price_col],
                "match_id":    match_id,
            })
            next_start_pos = i + 3   # AFTER MATCH SKIP PAST LAST ROW
    return pd.DataFrame(out_rows)

result = detect_two_consecutive_increases(df)
print(result)
\end{lstlisting}

\noindent\textbf{Results (Python/pandas):}
\begin{lstlisting}[basicstyle=\ttfamily\small, backgroundcolor=\color{gray!10}]
  company  start_price  end_price  match_id
0       A           36         36         1
1       A           30         30         2
2       B           24         24         1
\end{lstlisting}

For this reference scenario, the proposed engine, Trino~473, and the
hand-written Pandas baseline produce identical output.  This provides
evidence that partitioning, ordering, implicit \texttt{DEFINE}
behavior, navigation with \texttt{PREV}, \texttt{MATCH\_NUMBER()}, and
measure materialization are handled correctly for this supported
pattern family.

% ----------------------------------------------------------------
\section{Performance Results}
\label{sec:eval-performance}
% ----------------------------------------------------------------

Performance was evaluated on the Amazon UK Products Dataset using
five pattern families and five input sizes (25\,K, 35\,K, 50\,K,
75\,K, and 100\,K rows).  The detailed execution-time, throughput,
and memory results are reported in
Sections~\ref{subsec:eval-real-dataset}--\ref{subsec:eval-key-findings}.
Across the 25 measured runs, the engine completed all tests
successfully, with an average throughput of approximately
9\,838~rows/s and a total measured execution time of about
157.45\,s.  The results show predictable scaling over the evaluated
range: simpler patterns remain faster, while quantified and nested
patterns incur higher transition and condition-evaluation overhead.

% ----------------------------------------------------------------
\section{Comparison with Existing Systems}
\label{sec:eval-comparison}
% ----------------------------------------------------------------

The evaluated system is not intended to replace full SQL engines such
as Oracle or Trino, nor streaming platforms such as Flink.  Its design
goal is narrower: it brings a practical subset of SQL row-pattern
matching into in-memory Python analytics workflows over
\texttt{pandas} DataFrames.

Compared with Trino and Oracle, the proposed implementation avoids the
need to load data into a database server before running
\texttt{MATCH\_RECOGNIZE}.  This makes it convenient for notebook and
DataFrame-centered analysis, but it also means that the implementation
supports only the subset described in Chapters~\ref{chap:methodology}
and~\ref{chap:implementation}.  Trino and Oracle remain more complete
SQL systems with mature optimizers, distributed execution, and broader
standard coverage.  Compared with hand-written Pandas logic, the SQL
form is more declarative and reusable: partitioning, ordering, pattern
variables, skip behavior, and measures are expressed in one query
rather than in custom procedural code for each pattern.

% ----------------------------------------------------------------
\section{Use Cases and Applications}
\label{sec:eval-usecases}
% ----------------------------------------------------------------

The \texttt{MATCH\_RECOGNIZE} implementation supports a wide range of
applications in various fields, such as log analysis, time series
processing, and identifying patterns in financial sequences where
pattern recognition in sequential data is critical.

To assess correctness, each pattern-matching query was executed using
three independent approaches.  First, a \textbf{Python-UDF baseline}:
custom functions were implemented---built from \texttt{groupby},
\texttt{rolling}, and \texttt{shift}---to detect the same sequences
directly in \texttt{pandas} DataFrames; this required adjusting the
code for each query to reproduce the intended
semantics~\cite{Python_UDF}.  Second, a \textbf{Trino baseline}: the
equivalent \texttt{MATCH\_RECOGNIZE} statement was run in Trino
(version~473) and treated its output as the reference result.  Third,
the \textbf{proposed engine}: \texttt{MATCH\_RECOGNIZE} was executed
on the same SQL string and DataFrame; the system parses the query,
compiles a DFA, and performs the matching within \texttt{pandas}.

\noindent Throughout the evaluated scenarios, the results from the
three approaches were compared on a row-by-row basis.  For the
reference cases reported in this chapter, the
\texttt{MATCH\_RECOGNIZE} output matched the Trino result
(see Table~\ref{tab:sql-match-recognize-examples}).  In contrast, the
UDF method involved significantly more code and was more prone to
errors during development~\cite{Python_UDF}.

\begin{longtable}{|p{0.95\textwidth}|}
\caption{SQL \texttt{MATCH\_RECOGNIZE} Pattern Examples~\cite{monierashraf2025rowmatch}}
\label{tab:sql-match-recognize-examples}\\
\hline
\multicolumn{1}{|l|}{\parbox{0.93\textwidth}{\textbf{Example 1: Detecting `Dip' Patterns
(Three Orders with Successively Lower Prices)}}} \\
\hline
\textbf{SQL Query} \\
\begin{minipage}{0.90\textwidth}
\begin{verbatim}
SELECT customer_id, start_date, end_date, bottom_price
FROM orders
MATCH_RECOGNIZE (
  PARTITION BY customer_id
  ORDER BY order_date
  MEASURES
    A.order_date AS start_date,
    C.order_date AS end_date,
    C.price      AS bottom_price
  PATTERN (A B C)
  DEFINE
    B AS B.price < A.price,
    C AS C.price < B.price
);
\end{verbatim}
\end{minipage} \\
\hline
\textbf{Trino Results:} \\
\begin{minipage}{0.90\textwidth}
\begin{verbatim}
customer_id | start_date | end_date   | bottom_price
------------+------------+------------+--------------
cust_1      | 2020-05-12 | 2020-05-16 | 50
(1 row)
\end{verbatim}
\end{minipage} \\
\hline
\textbf{Proposed \texttt{MATCH\_RECOGNIZE} Results:} \\
\begin{minipage}{0.90\textwidth}
\begin{verbatim}
  customer_id  start_date    end_date  bottom_price
0      cust_1  2020-05-12  2020-05-16            50
\end{verbatim}
\end{minipage} \\
\hline
\textbf{Python UDF Results:} \\
\begin{minipage}{0.90\textwidth}
\begin{verbatim}
  customer_id  start_date    end_date  bottom_price
0      cust_1  2020-05-12  2020-05-16            50
\end{verbatim}
\end{minipage} \\
\hline
\end{longtable}

% ................................................................
\subsection{Performance Results Using Real Dataset}
\label{subsec:eval-real-dataset}
% ................................................................

The Amazon UK Products Dataset is a large e-commerce dataset
containing information about products listed on Amazon's UK
marketplace~\cite{asaniczka2023amazon}, commonly used for data
analysis, machine learning, and benchmarking.  It is provided as a
single CSV file with roughly 2.2\,million product records.  Each
record includes: \texttt{asin} (Amazon Standard Identification Number;
unique product ID); \texttt{title} (product name); \texttt{imgUrl}
(primary product image URL); \texttt{productURL} (canonical product
page URL); \texttt{stars} (average customer rating, 1--5,
floating-point); \texttt{reviews} (number of customer reviews,
non-negative integer); \texttt{price} (listed price in GBP, numeric);
\texttt{isBestSeller} (bestseller indicator, boolean);
\texttt{boughtInLastMonth} (approximate purchases in the last 30 days,
integer); and \texttt{categoryName} (top-level category label, e.g.\
Electronics, Books, Clothing).  This dataset is representative for
price-trend analysis, brand/feature-popularity studies, pattern
recognition, and recommender-system benchmarking.

The methodology involved testing 25 test cases across 5 SQL patterns
and 5 input sizes (25\,K, 35\,K, 50\,K, 75\,K, 100\,K rows).  All
tests completed successfully.  The implementation achieved an average
throughput of 9\,838~rows/s, and the end-to-end runtime over all 25
test cases was approximately 157.45\,s.  The results indicate
near-linear scaling over the evaluated range, with the expected
performance differences between simple, alternation, quantified,
optional, and nested patterns.

% ................................................................
\subsection{Pattern Definitions}
\label{subsec:eval-pattern-defs}
% ................................................................

{\small
\setlength{\LTleft}{0pt}
\setlength{\LTright}{0pt}
\begin{longtable}{@{\extracolsep{\fill}}llp{0.55\textwidth}}
\caption{SQL Pattern Definitions and Detection Goals}
\label{tab:pattern_definitions} \\
\toprule
\textbf{Pattern Name} & \textbf{SQL Pattern} & \textbf{Description \& Goal} \\
\midrule
\endfirsthead
\toprule
\textbf{Pattern Name} & \textbf{SQL Pattern} & \textbf{Description \& Goal} \\
\midrule
\endhead
\bottomrule
\endfoot

simple\_sequence & \texttt{A+ B+} &
\textbf{Goal:} Detect transitions from unrated/poor to very poor products.
\newline
\textbf{Use:} Identify sequences where no-rating/poor products (0--1.0
stars) are followed by very poor products (1.1--2.0 stars), useful for
detecting problematic product sequences or data quality issues in
listings. \\
\midrule

alternation & \texttt{A (B|C)+ D} &
\textbf{Goal:} Detect quality improvement patterns.
\newline
\textbf{Use:} Find sequences starting with unrated/poor products,
followed by mixed very poor/below-average products, ending with
average-to-good products.  Useful for analyzing quality improvement
patterns in product browsing sequences or sorted lists. \\
\midrule

quantified & \texttt{A\{2,5\} B* C+} &
\textbf{Goal:} Detect constrained low-quality patterns.
\newline
\textbf{Use:} Find sequences with 2--5 unrated/poor products,
optionally followed by very poor products, then below-average products.
Enforces minimum unrated product counts for identifying problematic
product clusters. \\
\midrule

optional\_pattern & \texttt{A+ B? C*} &
\textbf{Goal:} Flexible low-to-moderate quality transition detection.
\newline
\textbf{Use:} Broad pattern matching starting with unrated/poor
products with optional transitions through worse ratings.  High recall
for various low-quality sequence scenarios in e-commerce data. \\
\midrule

complex\_nested & \texttt{(A|B)+ (C\{1,3\} D*)+} &
\textbf{Goal:} Complex quality improvement analysis.
\newline
\textbf{Use:} Detect sequences of unrated/very poor products followed
by groups of improving quality (below-average to good).  Useful for
multi-level quality grouping and advanced sorting pattern detection in
product listings. \\
\end{longtable}
}

\noindent
All patterns analyze product quality based on star ratings where
$A$~= No Rating/Poor (0--1.0 stars),
$B$~= Very Poor (1.1--2.0 stars),
$C$~= Below Average (2.1--3.0 stars),
$D$~= Average to Good (3.1--4.0 stars),
$E$~= Excellent (4.1--5.0 stars).
Categories were created using equal-width binning on the \texttt{stars}
column.  Table~\ref{tab:pattern_definitions} provides detailed
specifications of all five SQL patterns used in the evaluation,
including their formal syntax and intended detection goals.  Each
pattern represents a different complexity level and use case, ranging
from simple sequential matching to complex nested structures with
quantifiers and alternation operators.

\begin{longtable}{lrcr}
\caption{Pattern Performance Summary (5 Patterns $\times$ 5 Dataset Sizes = 25 Tests)}
\label{tab:pattern_summary} \\
\toprule
\textbf{Pattern} & \textbf{Avg Throughput} & \textbf{Avg Time} & \textbf{Test Cases} \\
 & \textbf{(rows/s)} & \textbf{(s)} & \\
\midrule
\endfirsthead
\caption*{Table~\thetable{} (continued): Pattern Performance Summary} \\
\toprule
\textbf{Pattern} & \textbf{Avg Throughput} & \textbf{Avg Time} & \textbf{Test Cases} \\
 & \textbf{(rows/s)} & \textbf{(s)} & \\
\midrule
\endhead
\bottomrule
\endfoot
simple\_sequence  & 12{,}618 & 4.57 & 25K, 35K, 50K, 75K, 100K \\
alternation       & 10{,}881 & 5.34 & 25K, 35K, 50K, 75K, 100K \\
quantified        &  7{,}034 & 8.13 & 25K, 35K, 50K, 75K, 100K \\
optional\_pattern & 11{,}931 & 4.86 & 25K, 35K, 50K, 75K, 100K \\
complex\_nested   &  6{,}724 & 8.59 & 25K, 35K, 50K, 75K, 100K \\
\end{longtable}

Table~\ref{tab:pattern_summary} summarizes the performance
characteristics of each pattern averaged across all five dataset
sizes.  A clear inverse relationship exists between pattern complexity
and throughput: simple patterns achieve nearly double the throughput of
complex patterns.  The \texttt{simple\_sequence} pattern delivered the
best throughput at 12\,618~rows/s.  The \texttt{complex\_nested}
pattern, while lower at 6\,724~rows/s, successfully handles intricate
nested structures.  The \texttt{alternation} pattern maintains good
throughput at 10\,881~rows/s, and \texttt{optional\_pattern} provides
strong performance at 11\,931~rows/s.

\begin{table}[H]
\centering
\caption{Performance Summary by Dataset Size}
\label{tab:size_summary}
\begin{tabular}{rrrr}
\toprule
\textbf{Dataset Size} & \textbf{Avg Throughput} & \textbf{Avg Time}
  & \textbf{Test Cases} \\
\textbf{(rows)} & \textbf{(rows/s)} & \textbf{(s)} & \\
\midrule
 25{,}000 & 10{,}250 &  2.62 & 5 patterns \\
 35{,}000 &  9{,}840 &  3.83 & 5 patterns \\
 50{,}000 & 10{,}091 &  5.35 & 5 patterns \\
 75{,}000 &  9{,}495 &  8.46 & 5 patterns \\
100{,}000 &  9{,}512 & 11.22 & 5 patterns \\
\bottomrule
\end{tabular}
\end{table}

Table~\ref{tab:size_summary} presents performance metrics organized by
dataset size, averaging results across all five patterns.  The table
shows throughput stability across scales, with the mean throughput
remaining between approximately 9\,495 and 10\,250~rows/s across the
five evaluated input sizes.

\begin{longtable}{lrrrrr}
\caption{Execution Time (seconds) by Pattern and Dataset Size}
\label{tab:execution_times} \\
\toprule
\textbf{Pattern} & \multicolumn{5}{c}{\textbf{Dataset Size (rows)}} \\
\cmidrule(lr){2-6}
 & \textbf{25K} & \textbf{35K} & \textbf{50K} & \textbf{75K} & \textbf{100K} \\
\midrule
\endfirsthead
\caption*{Table~\thetable{} (continued): Execution Time by Pattern and Dataset Size} \\
\toprule
\textbf{Pattern} & \multicolumn{5}{c}{\textbf{Dataset Size (rows)}} \\
\cmidrule(lr){2-6}
 & \textbf{25K} & \textbf{35K} & \textbf{50K} & \textbf{75K} & \textbf{100K} \\
\midrule
\endhead
\bottomrule
\endfoot
simple\_sequence  & 1.94 &  2.77 &  3.82 &  6.12 &  8.20 \\
alternation       & 2.19 &  3.19 &  4.41 &  7.18 &  9.75 \\
quantified        & 3.45 &  5.07 &  7.06 & 10.87 & 14.19 \\
optional\_pattern & 1.98 &  2.92 &  4.13 &  6.58 &  8.69 \\
complex\_nested   & 3.55 &  5.22 &  7.31 & 11.57 & 15.29 \\
\end{longtable}

Table~\ref{tab:execution_times} presents a matrix of execution times
for all pattern--size combinations.  The table illustrates an
approximately linear relationship between dataset size and execution
time over the evaluated range.  The \texttt{complex\_nested} pattern
consistently
requires the longest execution time (15.29\,s for 100\,K rows), while
\texttt{simple\_sequence} achieves the fastest processing
(1.94\,s for 25\,K rows).

% ................................................................
\subsection{Throughput by Pattern and Size}
\label{subsec:eval-throughput}
% ................................................................

\enlargethispage{2\baselineskip}

{\small
\begin{longtable}{@{\extracolsep{\fill}}lrrrrr}
\caption{Throughput (rows/s) by Pattern and Dataset Size}
\label{tab:throughput} \\
\toprule
\textbf{Pattern} & \multicolumn{5}{c}{\textbf{Dataset Size (rows)}} \\
\cmidrule(lr){2-6}
 & \textbf{25K} & \textbf{35K} & \textbf{50K} & \textbf{75K} & \textbf{100K} \\
\midrule
\endfirsthead
\midrule
\endhead
\bottomrule
\endfoot
simple\_sequence  & 12{,}918 & 12{,}619 & 13{,}097 & 12{,}256 & 12{,}202 \\
alternation       & 11{,}402 & 10{,}979 & 11{,}328 & 10{,}441 & 10{,}255 \\
quantified        &  7{,}243 &  6{,}901 &  7{,}082 &  6{,}898 &  7{,}048 \\
optional\_pattern & 12{,}642 & 11{,}993 & 12{,}108 & 11{,}400 & 11{,}513 \\
complex\_nested   &  7{,}045 &  6{,}710 &  6{,}842 &  6{,}481 &  6{,}542 \\
\end{longtable}
}

\noindent\textbf{Throughput (rows/s)} measures processing speed as
dataset size divided by execution time; higher values indicate faster
evaluation.  Across all experiments, throughput ranged from
6{,}481 to 13{,}097~rows/s (mean 9{,}838~rows/s).  Simple sequential
patterns, which require fewer state transitions, typically achieve
12{,}000--13{,}000~rows/s, whereas complex nested patterns with
multiple evaluation layers achieve about 6{,}000--7{,}000~rows/s.
For a fixed pattern, throughput remains relatively constant as input
size increases, supporting the near-linear scaling trend observed in
Table~\ref{tab:execution_times}.

Table~\ref{tab:throughput} displays throughput measurements across
all test scenarios.  The table shows that most variation is driven by
pattern complexity rather than input size: simple and optional
patterns remain fastest, while quantified and nested patterns are
consistently slower.

% ................................................................
\subsection{Memory Metrics}
\label{subsec:eval-memory}
% ................................................................

\noindent\textbf{Memory Usage (MB)} represents the average memory
consumed during pattern matching operations.  This includes memory
for storing intermediate matching states, row buffers, and result
sets.  Memory usage varies based on both dataset size and pattern
complexity.  Simple patterns with fewer matches may use less memory
(0.56--3.20\,MB for small datasets), while patterns that generate
large result sets require more memory (15--30\,MB).  The memory
values shown have been corrected to absolute values, as some
measurements showed negative artifacts due to garbage collection
timing.

\noindent\textbf{Peak Memory (MB)} indicates the maximum memory
consumption during pattern evaluation, typically occurring during
result collection phases.  Peak memory is approximately $1.3\times$
the average memory usage, accounting for temporary allocations during
pattern state transitions and match aggregation.  Across all 25 test
cases, peak memory remains under 43\,MB (maximum 42.34\,MB at 75\,K
rows), indicating modest memory use for the evaluated workloads rather
than a strictly monotonic relationship with dataset size.  Notably,
the 100\,K-row tests consume less peak memory
(maximum 39.15\,MB) than the 75\,K-row tests, further validating the
non-linear memory behavior.

Memory consumption does not correlate linearly with dataset size or
match count, but rather depends on the complexity of intermediate
pattern states and result-set characteristics.  For example, at
100\,K rows, \texttt{optional\_pattern} finds 15{,}247 matches
($8.4\times$ more than \texttt{alternation}'s 1{,}828) but uses only
1.36\,MB compared to \texttt{alternation}'s 30.12\,MB, showing
that match count alone does not determine memory usage.  Most
notably, \texttt{simple\_sequence} at 75\,K rows consumes
42.34\,MB peak memory, while the same pattern at 100\,K rows uses
only 30.38\,MB---a 28\% reduction despite the larger dataset.  This
variability suggests memory usage is driven more by pattern state
complexity, garbage collection timing, and internal representation
than by output size or dataset scale.  Despite these variations, the
implementation demonstrates efficient memory utilization, with peak
memory remaining within acceptable bounds (under 43\,MB) across all
test scenarios.

{\small
\setlength{\LTleft}{0pt}
\setlength{\LTright}{0pt}
\begin{longtable}{@{\extracolsep{\fill}}rlrrr}
\caption{Memory Consumption Metrics}
\label{tab:memory_cache} \\
\toprule
\textbf{Dataset Size} & \textbf{Pattern} & \textbf{Execution} & \textbf{Memory} & \textbf{Peak Memory} \\
\textbf{(rows)} & \textbf{Complexity} & \textbf{Time (ms)} & \textbf{Usage (MB)} & \textbf{(MB)} \\
\midrule
\endfirsthead
\caption*{Table~\thetable{} (continued): Memory Consumption Metrics} \\
\toprule
\textbf{Dataset Size} & \textbf{Pattern} & \textbf{Execution} & \textbf{Memory} & \textbf{Peak Memory} \\
\textbf{(rows)} & \textbf{Complexity} & \textbf{Time (ms)} & \textbf{Usage (MB)} & \textbf{(MB)} \\
\midrule
\endhead
\bottomrule
\endfoot
 25{,}000 & simple\_sequence  & 1{,}935 & 15.20 & 19.76 \\
 25{,}000 & alternation       & 2{,}193 &  2.51 &  3.27 \\
 25{,}000 & optional\_pattern & 1{,}978 &  6.73 &  8.75 \\
 25{,}000 & quantified        & 3{,}451 &  1.02 &  1.33 \\
 25{,}000 & complex\_nested   & 3{,}548 &  0.56 &  0.73 \\
\midrule
 35{,}000 & simple\_sequence  & 2{,}774 & 13.21 & 17.17 \\
 35{,}000 & alternation       & 3{,}187 &  1.23 &  1.60 \\
 35{,}000 & optional\_pattern & 2{,}918 &  7.68 &  9.98 \\
 35{,}000 & quantified        & 5{,}073 &  3.68 &  4.78 \\
 35{,}000 & complex\_nested   & 5{,}216 & 10.95 & 14.24 \\
\midrule
 50{,}000 & simple\_sequence  & 3{,}817 & 11.27 & 14.65 \\
 50{,}000 & alternation       & 4{,}413 &  2.60 &  3.38 \\
 50{,}000 & optional\_pattern & 4{,}129 &  2.80 &  3.64 \\
 50{,}000 & quantified        & 7{,}059 &  4.66 &  6.06 \\
 50{,}000 & complex\_nested   & 7{,}306 &  6.57 &  8.54 \\
\midrule
 75{,}000 & simple\_sequence  & 6{,}120 & 32.57 & 42.34 \\
 75{,}000 & alternation       & 7{,}181 &  8.71 & 11.32 \\
 75{,}000 & optional\_pattern & 6{,}577 &  8.86 & 11.52 \\
 75{,}000 & quantified        & 10{,}872 &  5.93 &  7.71 \\
 75{,}000 & complex\_nested   & 11{,}572 & 15.51 & 20.16 \\
\midrule
100{,}000 & simple\_sequence  &  8{,}195 & 23.37 & 30.38 \\
100{,}000 & alternation       &  9{,}750 & 30.12 & 39.15 \\
100{,}000 & optional\_pattern &  8{,}686 &  1.36 &  1.77 \\
100{,}000 & quantified        & 14{,}188 & 20.60 & 26.78 \\
100{,}000 & complex\_nested   & 15{,}286 & 17.82 & 23.17 \\
\end{longtable}
}

Table~\ref{tab:memory_cache} presents detailed memory consumption
metrics for all test scenarios, tracking both average and peak memory
usage.  The table shows bounded memory usage for the evaluated
workloads, with peak memory remaining under 43\,MB across all 25 test
cases.  The maximum peak memory of 42.34\,MB occurs at 75\,K
rows (\texttt{simple\_sequence}) rather than at the largest dataset
size (100\,K rows), illustrating that memory consumption is
non-linear and depends on runtime factors including garbage
collection timing, pattern state complexity, and the specific
distribution of category transitions in the data.

% ................................................................
\subsection{Key Findings}
\label{subsec:eval-key-findings}
% ................................................................

\begin{table}[H]
\centering
\caption{Overall Test Statistics}
\label{tab:overall_stats}
\begin{tabular}{lr}
\toprule
\textbf{Metric} & \textbf{Value} \\
\midrule
Total Tests             & 25 \\
Success Rate            & 100\% \\
Total Execution Time    & 157.45\,s \\
Average Execution Time  & 6.30\,s \\
Average Throughput      & 9{,}838 rows/s \\
Min Throughput          & 6{,}481 rows/s \\
Max Throughput          & 13{,}097 rows/s \\
\bottomrule
\end{tabular}
\end{table}

Table~\ref{tab:overall_stats} presents the comprehensive test
statistics aggregated across all 25 test cases.  The system completed
all evaluated patterns and dataset sizes successfully, with no
failures or errors encountered during the entire evaluation suite.
Linear scaling is demonstrated from 25\,K to 100\,K rows, with
execution time increasing proportionally and predictably as dataset
size grows.  The implementation successfully handles patterns ranging
from simple sequences to complex nested structures, providing evidence
of practical viability for medium-scale sequential pattern detection
over e-commerce data.

% ----------------------------------------------------------------
\section{Chapter Summary}
\label{sec:eval-summary}
% ----------------------------------------------------------------

This chapter evaluated the proposed \texttt{MATCH\_RECOGNIZE}
implementation through correctness checks, feature-coverage tests,
performance measurements, memory analysis, and comparison with
existing systems.  The experiments used five pattern families across
five dataset sizes, producing 25 successful performance runs.

The results show that the implementation preserves the evaluated
semantics while reducing developer effort compared with hand-written
Pandas logic.  The next chapter summarizes the thesis contributions,
discusses limitations, and outlines future work.
